\subsection{局部紧群的星代数}

定义 9.　设 G 为局部紧群，A 为 G 上关于 G 的某个 Haar 测度具有密度的有界测度所成的对合 Banach 代数（I，第 99 页，例 4）。对合 Banach 代数 A 的包络星代数称为 G 的星代数，记作 Stell(G)。

注.　设 \(\nu\) 为 G 上的一个左 Haar 测度，\(\Delta\) 为其模函数。映射 \(f \mapsto f \cdot \nu\) 是代数 \(\mathrm{L}^{1}(\mathrm{G}, \nu)\) 到 A 上的等距同构（同前）。因此也可以把 Stell(G) 定义为赋范对合代数 \(\mathrm{L}^{1}(\mathrm{G}, \nu)\) 的包络星代数。

设 G 为局部紧群，\(\nu\) 为 G 上的一个左 Haar 测度。对 \(\mu \in \mathcal{M}^1 (\mathrm{G})\) 与 \(f\in \mathrm{L}^2 (\mathrm{G},\nu)\) ，有 \(\mu *f\in \mathrm{L}^2 (\mathrm{G},\nu)\) （INT，VIII，§4，命题 6）。于是，我们用 \(\pmb{\gamma}(\mu)\) 表示由 \(f\mapsto \mu *f\) 所给出的自同态，它作用于 \(\mathrm{L}^2 (\mathrm{G},\nu)\) 。映射 \(\mu \mapsto \pmb{\gamma}(\mu)\) 给出代数 \(\mathcal{M}^{1}(\mathrm{G})\) 在 Banach 代数 \(\mathcal{L}(\mathrm{L}^{2}(\mathrm{G},\nu))\) （由 \(\mathrm{L}^2 (\mathrm{G},\nu)\) 的连续自同态组成）中的一个表示（INT，VIII，§4，命题 6 的推论）。另一方面，\(\gamma (\breve{\mu})\) 是自同态 \(\gamma (\mu)\) 的转置（INT，VIII，§4，第 3 节，命题 8）。由此得出 \(\gamma (\mu^{*})\) 是 \(\gamma (\mu)\) 的伴随自同态，从而映射 \(\gamma \colon \mu \mapsto \gamma (\mu)\) 是从 \(\mathcal{M}^{1}(\mathrm{G})\) 到星代数 \(\mathcal{L}(\mathrm{L}^{2}(\mathrm{G},\nu))\) 中的对合代数态射，称为 \(\mathcal{M}^{1}(\mathrm{G})\) 在 \(\mathrm{L}^2 (\mathrm{G},\nu)\) 中的左正则表示。根据 INT，VIII，§4，第 7 节，命题 19，这个表示是忠实的。

设 \(j\) 为从 \(\mathrm{L}^{1}(\mathrm{G},\nu)\) 到 \(\mathrm{Stell}(\mathrm{G})\) 中的典范映射。通过限制到 \(\mathrm{L}^{1}(\mathrm{G},\nu)\) 上，正则表示 \(\pmb{\gamma}\) 便给出从 \(\mathrm{L}^{1}(\mathrm{G},\nu)\) 到 \(\mathcal{L}(\mathrm{L}^{2}(\mathrm{G},\nu))\) 中的一个单射对合代数态射，称为 \(\mathrm{L}^{1}(\mathrm{G})\) 在 \(\mathrm{L}^{2}(\mathrm{G})\) 中的左正则表示。根据命题 20，存在唯一的态射 \(\pmb{\gamma}'\colon \mathrm{Stell}(\mathrm{G})\to \mathcal{L}(\mathrm{L}^{2}(\mathrm{G},\nu))\) 使得 \(\pmb{\gamma} = \pmb{\gamma}'\circ j\) 。我们称 \(\pmb{\gamma}'\) 为 \(\mathrm{Stell}(\mathrm{G})\) 在 \(\mathrm{L}^{2}(\mathrm{G},\nu)\) 中的左正则表示。为记号简便起见，我们仍把它记作

\[
\boldsymbol {\gamma} ^ {\prime} (\varphi) (f) = \varphi * f\tag{7}
\]

其中 \(f\in \mathrm{L}^2 (\mathrm{G},\nu)\) ，\(\varphi \in \mathrm{Stell}(\mathrm{G})\) 。我们有

\[
\| \varphi * f \| _ {2} \leqslant \| \varphi \| _ {*} \| f \| _ {2}.\tag{8}
\]

注.　一般说来，Stell(G) 的左正则表示 \(\gamma'\) 在 \(L^{2}(G,\nu)\) 中不是忠实的。可以证明，它为忠实表示的充分必要条件是：\(\mathrm{L}_{\mathbf{R}}^{\infty}(\mathrm{G})\) 上存在正线性型 f，满足 \(f(1)=1\) ，且 \(f(\boldsymbol{\gamma}(g)x)=f(x)\) 对所有 \((g,x)\in\mathrm{G}\times\mathrm{G}\) 成立（这时称群 G 为顺从群，参见 EVT，IV，第 73 页，习题 4）。

命题 22.　从 \(\mathrm{L}^{1}(\mathrm{G},\nu)\) 到 Stell(G) 中的典范映射 j 是单射，且像稠密。

根据赋范对合代数的包络星代数的定义，\(j\) 的像是稠密的。由于左正则表示 \(\gamma\) 是忠实的，故 \(j\) 的单射性由等式 \(\gamma = \gamma' \circ j\) 得出。

推论.　代数 \(\mathrm{L}^{1}(\mathrm{G},\nu)\) 无根。

这由 I，第 125 页，命题 21 与命题 22 得出。

于是可以把 \(\mathrm{L}^{1}(\mathrm{G},\nu)\) 等同于 \(\mathrm{Stell}(\mathrm{G})\) 的一个稠密对合子代数，而 \(\mathrm{L}^{1}(\mathrm{G},\nu)\) 到 \(\mathrm{Stell}(\mathrm{G})\) 中的典范单射这时是连续的。

现在假设群 G 是幺模群（INT，VII，§1，第 3 节，定义 3）。于是可以从右正则表示 \((f,\mu)\mapsto\boldsymbol{\delta}(\mu)(f)=f*\breve{\mu}\) （从 \(\mathrm{L}^{2}(\mathrm{G},\nu)\times\mathcal{M}^{1}(\mathrm{G})\) 到 \(\mathrm{L}^{2}(\mathrm{G},\nu)\) 中的表示）出发重复同样的论证。这样就定义了一个态射 \(\delta^{\prime}\) ，它从 \(\operatorname{Stell}(\mathrm{G})\) 到 \(\mathcal{L}(\mathrm{L}^{2}(\mathrm{G},\nu))\) ，满足 \(\delta=\delta^{\prime}\circ j\) ，并定义 \(\boldsymbol{\delta}^{\prime}(\varphi)(f)=f*\varphi\) ，其中 \(f\in\mathrm{L}^{2}(\mathrm{G},\nu)\) ，\(\varphi\in\operatorname{Stell}(\mathrm{G})\) 。

对 \(\varphi, \psi \in \text{Stell}(G)\) ，有 \(\boldsymbol{\delta}'(\psi) \circ \boldsymbol{\gamma}'(\varphi) = \boldsymbol{\gamma}'(\psi) \circ \boldsymbol{\delta}'(\varphi)\) ，即

\[
(\varphi * f) * \psi = \varphi * (f * \psi)\tag{9}
\]

对所有 \(f \in \mathrm{L}^{2}(\mathrm{G}, \nu)\) 成立。事实上，该公式当 \(\varphi, \psi \in \mathrm{L}^{1}(\mathrm{G}, \nu)\) 时成立，且映射 \((\varphi, \psi) \mapsto (\varphi * f) * \psi\) 与 \((\varphi, \psi) \mapsto \varphi * (f * \psi)\) 都是从 \(\operatorname{Stell}(\mathrm{G}) \times \operatorname{Stell}(\mathrm{G})\) 到 \(\mathrm{L}^{2}(\mathrm{G}, \nu)\) 中的连续双线性映射。

